\subsection{Finite central recovery in the mixed-braid target}

\begin{result}{A.4.28}{Proved}{The mixed coefficient test has an exact finite coloured-orbifold target}
\label{prop:mixed-finite-central-recovery}
Put
\[
\begin{gathered}
 H=\langle u,v\mid(uv)^2=(vu)^2\rangle,
 \qquad z=(uv)^2,\\
 R_0=H/\langle z\rangle
       =\langle u,v\mid(uv)^2=1\rangle\cong C_2*\mathbb Z.
\end{gathered}                                                  \tag{MCR.1}
\]
For every finite image $q:H\twoheadrightarrow Q$, define
$n=\operatorname{ord}(q(z))$ and the actual homomorphisms
\[
\begin{gathered}
 \chi:H\longrightarrow\mathbb Z/(2n),\qquad
       \chi(u)=1,\quad\chi(v)=0,\\
 \widetilde q=(q,\chi),\qquad
 \widetilde Q=\widetilde q(H),\qquad
 \zeta=\widetilde q(z),\\
 C_q=\widetilde Q/\langle\zeta\rangle,\qquad
 \pi_q:R_0\twoheadrightarrow C_q.
\end{gathered}                                                  \tag{MCR.2}
\]
Then $|C_q|\leq2|Q|$.  For every commutator proword
$f\in[\widehat F_2,\widehat F_2]$ and $c\in\widehat{\mathbb Z}$,
the following two finite identities are equivalent:
\[
\begin{aligned}
 f(q(u),q(v)^2)
   &=q(v)^{4c}f(q(u)^2,q(v))q(u)^{4c}q(z)^{-2c},\\
 f(\pi_q(u),\pi_q(v)^2)
   &=\pi_q(v)^{4c}f(\pi_q(u)^2,\pi_q(v))\pi_q(u)^{4c}.
\end{aligned}                                                   \tag{MCR.3}
\]
Thus the first identity, whose $(u,v)=(d,V)$ specialization is
(ABR.4), is recovered from the displayed finite quotient of the
actual mixed four-point orbifold group, with its central coordinate
retained by the map (MCR.2).

Let $S:R_0\to R_0$ interchange $u,v$.  The joint image
\[
 C_q^{\mathrm{sym}}
    =\operatorname{im}\bigl((\pi_q,\pi_q S):
                          R_0\longrightarrow C_q\times C_q\bigr)
                                                               \tag{MCR.4}
\]
has both indicated quotient projections and the automorphism
$(x,y)\mapsto(y,x)$.  For the actual mixed image on the faithful
quaternion module together with the signed compact periods modulo
four, the orders and generator orders are
\[
\begin{array}{c|r|c}
\text{group}&\text{order}&\text{orders of the marked }u,v\\\hline
 Q&432&(4,4)\\
 \widetilde Q&1296&(12,4)\\
 C_q&216&(12,4)\\
 C_q^{\mathrm{sym}}&648&(12,12).
\end{array}                                                     \tag{MCR.5}
\]
Here $q(z)$ has order six, and the derived subgroup of $Q$ is
the nonabelian exponent-three group of order $27$.  Its image
in $C_q$ is still isomorphic to that same group.
\end{result}

\begin{proof}
The relation of $H$ makes $z$ central: multiplication by $u$
or $v$ on either side of $(uv)^2=(vu)^2$ gives the two required
commutations.  The presentation of $R_0$ in (MCR.1) is exact;
putting $s=uv$ gives the inverse generator changes
$s\mapsto uv$ and $u\mapsto sv^{-1}$ between that presentation
and $\langle s,v\mid s^2=1\rangle$.

Since $\chi(z)=2$, the element
$\zeta=(q(z),2)$ has order exactly $n$, and its projection to
$\mathbb Z/(2n)$ has that same order.  Consequently
\[
 \langle\zeta\rangle\cap\ker\bigl(
      \widetilde Q\longrightarrow\mathbb Z/(2n)\bigr)=1.
                                                               \tag{MCR.6}
\]
The order bound follows from
$|\widetilde Q|\leq2n|Q|$ and $|\langle\zeta\rangle|=n$.
Killing $\zeta$ kills the image of $z$, giving precisely the
surjection $\pi_q$ from $R_0$.

Write the mixed defect in $\widehat H$ as
\[
 \mathcal D_c(f)=f(u,v^2)
       \bigl(v^{4c}f(u^2,v)u^{4c}z^{-2c}\bigr)^{-1}.
\]
Its two coordinates in
$\widehat{H^{\mathrm{ab}}}=\widehat{\mathbb Z}[u]
                       \oplus\widehat{\mathbb Z}[v]$ are zero:
the two evaluations of $f$ contribute zero, and the remaining
coordinates on the right are
$(4c,4c)-2c(2,2)=(0,0)$.
Thus $\chi(\mathcal D_c(f))=0$.  Vanishing of its image in
$C_q$ puts $\widetilde q(\mathcal D_c(f))$ in
$\langle\zeta\rangle$, and (MCR.6) then makes it trivial.
This proves the implication to the first identity in (MCR.3).
Conversely that first identity and the zero $\chi$ coordinate
make $\widetilde q(\mathcal D_c(f))=1$, which proves the second.
All maps to the finite groups extend continuously to full
completions, so the argument applies to the stated prowords
and profinite exponents without a choice of lifts.

The map $S$ is an involutive automorphism because it carries
$(uv)^2$ to its conjugate $(vu)^2$.  Applying $S$ before
$(\pi_q,\pi_q S)$ interchanges the two coordinates.  The image
in (MCR.4) is therefore preserved by the coordinate swap,
and both projections are surjective.  The test in (MCR.3)
is the test on its first quotient projection; retaining both
projections does not change that test.

For the numerical statement the two input matrices are the
literal local chain images $D=q(d)$ and $L=q(V)$, with
$V=(aba)^2$, on the faithful $36$-dimensional quaternion
homology and on the permutation representation of the sixteen
compact homology vectors modulo four.  This gives a faithful
$52$-dimensional signed matrix model over $\mathbb F_3$.
The independent program
\path{work/ghost_code/check_main_mixed_capping.py}
starts from these two matrices and rebuilds the entire image
by multiplication, retaining a positive source word for each
of its $432$ elements.  It verifies that $(DL)^2$ is central
of order six.  It then performs the construction (MCR.2)
literally as pairs $(g,j)$ with $j\in\mathbb Z/12$, using
generators $(D,1)$ and $(L,0)$.  There are $1296$ such pairs,
with three lifts of every original matrix.

The program partitions those pairs into the six-element
orbits of multiplication by $((DL)^2,2)$, giving $216$
cosets.  It checks that both generator transitions descend,
and that $(uv)^2$ acts trivially on every coset.  The two
marked orders are $12$ and $4$.  Breadth-first enumeration
of the paired generators $(\pi_q(u),\pi_q(v))$ and
$(\pi_q(v),\pi_q(u))$ gives the $648$-element image (MCR.4);
all its transitions intertwine with the literal coordinate
swap.  Both projections have all $216$ elements in their
image.

Finally the independently rebuilt full multiplication table
gives $|Q'|=27$, $|[Q',Q']|=3$, exponent three for $Q'$,
and centrality of its commutator.  Every $g\in Q'$ occurs
as $(g,0)$ in the joined group.  The quotient map takes these
to $27$ distinct cosets, as also follows from (MCR.6).
This proves the last assertion and records the retained
central order-three information explicitly.  The actual
text, JSON and matrix certificates include every source
word, every joined coordinate, the central cosets, the
generator transitions and both projections.

The equivalence (MCR.3) isolates the finite target of the
coefficient identity.  Its derivation from a finite pentagon
input remains the separate comparison required in A.4.
The present result imposes no extra equation on a pentagon
shadow and makes no assertion that such a shadow satisfies
either side of (MCR.3).
\end{proof}
